% ============================================================
% Heart Failure Survival Analysis — Research Poster
% beamerposter · landscape 120×85 cm · full-page fill
% ============================================================
\documentclass[final]{beamer}

\usepackage[
  orientation=landscape,
  size=custom,
  width=120,
  height=85,
  scale=1.1
]{beamerposter}

\usepackage[T1]{fontenc}
\usepackage{lmodern}
\usepackage{booktabs}
\usepackage{array}
\usepackage{tikz}
\usetikzlibrary{positioning,arrows.meta,calc}
\usepackage[table]{xcolor}
\usepackage{tcolorbox}
\tcbuselibrary{skins}
\usepackage{enumitem}
\usepackage{amsmath}
\usepackage{tabularx}
\usepackage{ragged2e}

% ── Palette ──────────────────────────────────────────────
\definecolor{Navy}{HTML}{1e3a5f}
\definecolor{Blue}{HTML}{2563eb}
\definecolor{LBlue}{HTML}{dbeafe}
\definecolor{Green}{HTML}{15803d}
\definecolor{LGreen}{HTML}{dcfce7}
\definecolor{Amber}{HTML}{b45309}
\definecolor{LAmber}{HTML}{fef3c7}
\definecolor{TGray}{HTML}{334155}
\definecolor{LGray}{HTML}{64748b}
\definecolor{Border}{HTML}{cbd5e1}

% ── Strip beamer chrome ──────────────────────────────────
\setbeamertemplate{navigation symbols}{}
\setbeamertemplate{headline}{}
\setbeamertemplate{footline}{}
\setbeamercolor{background canvas}{bg=white}

% ── Block style ──────────────────────────────────────────
\setbeamercolor{block title}{fg=white, bg=Navy}
\setbeamercolor{block body}{fg=TGray, bg=white}
\setbeamerfont{block title}{size=\large, series=\bfseries}
\setbeamerfont{block body}{size=\normalsize}

\setbeamertemplate{block begin}{%
  \vskip0.5cm
  \begin{beamercolorbox}[
    wd=\linewidth, ht=3.2ex, dp=1.0ex,
    leftskip=0.8em, rightskip=0.5em
  ]{block title}%
    \usebeamerfont{block title}\insertblocktitle
  \end{beamercolorbox}%
  \begin{beamercolorbox}[
    wd=\linewidth,
    leftskip=0.8em, rightskip=0.8em,
    sep=0.7em
  ]{block body}%
    \usebeamerfont{block body}%
}
\setbeamertemplate{block end}{%
  \end{beamercolorbox}%
  \vskip0.15cm%
}

% ── Finding cards ────────────────────────────────────────
\tcbset{
  findcard/.style={
    enhanced, sharp corners,
    colback=#1!8, colframe=#1!50,
    boxrule=1.5pt,
    left=7pt, right=7pt, top=5pt, bottom=5pt,
    fontupper=\normalsize,
  }
}

\setlength{\columnsep}{0.9cm}
\setlist[itemize]{leftmargin=1.3em, itemsep=3pt, topsep=3pt}

% ============================================================
\begin{document}
\begin{frame}[t]{}

% ── HEADER ───────────────────────────────────────────────
\begin{tikzpicture}[remember picture, overlay]
  \fill[Navy] (0,-5.8cm) rectangle (\paperwidth,0);
  \node[anchor=west, text=white,
        font=\fontsize{52}{60}\selectfont\bfseries]
    at (2cm,-1.3cm) {Heart Failure Survival Analysis};
  \node[anchor=west, text=LBlue, font=\Large]
    at (2cm,-3.0cm)
    {Machine Learning Prediction from Clinical Records
     \enspace\textbullet\enspace
     Michigan Data Science Team, Winter 2026};
  \node[anchor=west, text=Border, font=\large]
    at (2cm,-4.6cm)
    {Chicco \& Jurman (2020)\ \emph{BMC Medical Informatics
     and Decision Making}\ 20(16)};
\end{tikzpicture}
\vspace{6.3cm}

% ── BODY ─────────────────────────────────────────────────
\hspace{0.4cm}%
\begin{columns}[t, totalwidth=119.2cm]

%=============================================================
% COL 1 — Background · Dataset · Statistics · Unsupervised
%=============================================================
\begin{column}{28.8cm}

\begin{block}{Background \& Motivation}
  Heart failure is a leading global cause of mortality. This
  project applies a complete 9-week ML curriculum to \textbf{299}
  clinical patient records to determine whether
  \textbf{routine biomarkers} alone can reliably predict survival.

  \medskip
  \textbf{Core question:} Can we match full-feature model accuracy
  using only the 2 most informative biomarkers?

  \medskip
  \textbf{Approach:} Statistical testing, unsupervised clustering,
  supervised classification, ensemble methods, feature selection,
  deep learning, and SHAP explainability — applied end-to-end on
  the same dataset across 9 weeks.
\end{block}

\begin{block}{Dataset — UCI Heart Failure Records}
  \medskip
  \begin{tikzpicture}
    \pgfmathsetmacro{\survW}{26*0.679}
    \pgfmathsetmacro{\midS}{\survW/2}
    \pgfmathsetmacro{\midD}{\survW + (26-\survW)/2}
    \fill[Blue!65]  (0,0)         rectangle (\survW, 0.75);
    \fill[Amber!70] (\survW,0)    rectangle (26,     0.75);
    \node[white, font=\normalsize\bfseries, anchor=center]
      at (\midS, 0.375) {Survived\ \ 203\ \ (67.9\%)};
    \node[white, font=\normalsize\bfseries, anchor=center]
      at (\midD, 0.375) {Deceased\ \ 96\ \ (32.1\%)};
  \end{tikzpicture}

  \medskip
  \setlength{\tabcolsep}{6pt}
  \renewcommand{\arraystretch}{1.35}
  \begin{tabularx}{\linewidth}{@{}Xcc@{}}
    \toprule
    \textbf{Feature} & \textbf{Type} & \textbf{Range} \\
    \midrule
    \rowcolor{LBlue}\textbf{Ejection Fraction}$^\star$ & Numeric & 14–80\% \\
    \rowcolor{LBlue}\textbf{Serum Creatinine}$^\star$  & Numeric & 0.5–9.4 mg/dL \\
    Age                & Numeric & 40–95 yr \\
    Follow-up Time     & Numeric & 4–285 days \\
    Serum Sodium       & Numeric & 114–148 mEq/L \\
    CPK, Platelets     & Numeric & variable \\
    Anaemia / Diabetes & Binary  & 0 / 1 \\
    High BP / Smoking  & Binary  & 0 / 1 \\
    \textbf{DEATH\_EVENT} & Target & 0 / 1 \\
    \bottomrule
  \end{tabularx}
  \smallskip
  {\small $^\star$ Top predictors confirmed by MRMR, Lasso, and SHAP}
\end{block}

\begin{block}{Statistical Analysis — Week 2}
  T-test + Mann-Whitney U with Benjamini-Hochberg FDR
  correction. Significant features (p\,<\,0.05):

  \medskip
  \setlength{\tabcolsep}{6pt}
  \renewcommand{\arraystretch}{1.35}
  \begin{tabularx}{\linewidth}{@{}Xr@{}}
    \toprule
    \textbf{Feature} & \textbf{p-value} \\
    \midrule
    \rowcolor{LGreen}Follow-up Time    & $2.3\times10^{-22}$ \\
    \rowcolor{LGreen}Ejection Fraction & $9.6\times10^{-6}$  \\
    \rowcolor{LGreen}Serum Creatinine  & $6.4\times10^{-5}$  \\
    \rowcolor{LGreen}Age               & $4.7\times10^{-5}$  \\
    Serum Sodium      & $<0.05$             \\
    Anaemia, Diabetes, Smoking, Sex & n.s. \\
    \bottomrule
  \end{tabularx}
\end{block}

\begin{block}{Unsupervised Learning — Week 3}
  \textbf{PCA:} 9 components needed for 80\% explained variance
  — no single dominant structure, motivating supervised methods.

  \medskip
  \textbf{K-Means} (k=2): 165 vs.\ 134 patients.
  Clusters show moderate alignment with outcome labels but
  cannot reliably separate survival classes alone.

  \medskip
  \textbf{Hierarchical clustering:} Agglomerative linkage with
  elbow method and silhouette analysis both confirm k=2 as the
  natural partition.
\end{block}

\end{column}

%=============================================================
% COL 2 — Pipeline · Feature Selection
%=============================================================
\begin{column}{28.8cm}

\begin{block}{9-Week ML Pipeline}
  \begin{center}
  \begin{tikzpicture}[
    nd/.style={
      rectangle, rounded corners=5pt,
      draw=Blue!40, fill=LBlue,
      minimum width=10.5cm, minimum height=0.85cm,
      font=\normalsize\bfseries, text=Navy, align=center
    },
    ar/.style={-Stealth, very thick, Blue!55},
    node distance=0.30cm
  ]
    \node[nd](w1){W1\enspace EDA \& Visualisation};
    \node[nd,below=of w1](w2){W2\enspace Hypothesis Testing (T-test, FDR)};
    \node[nd,below=of w2](w3){W3\enspace PCA \& Clustering};
    \node[nd,below=of w3](w4){W4\enspace Supervised Classifiers};
    \node[nd,below=of w4](w5){W5\enspace Hyperparameter Optimisation};
    \node[nd,below=of w5](w6){W6\enspace Ensemble Methods};
    \node[nd,below=of w6](w7){W7\enspace Feature Selection};
    \node[nd,below=of w7](w8){W8\enspace Deep Learning (MLP)};
    \node[nd,below=of w8](w9){W9\enspace SHAP \& PLS-DA};
    \foreach \s/\t in
      {w1/w2,w2/w3,w3/w4,w4/w5,w5/w6,w6/w7,w7/w8,w8/w9}
      {\draw[ar](\s)--(\t);}
  \end{tikzpicture}
  \end{center}

  \medskip
  \textbf{Train/Test split:} 70/30 stratified (209 train / 90 test).
  All scaling and feature selection fitted on training data only to
  prevent leakage. LightGBM tuned with Optuna Bayesian
  optimisation over 100 trials.
\end{block}

\begin{block}{Feature Selection — Week 7}
  Three independent methods converge on the same top features:

  \medskip
  \setlength{\tabcolsep}{6pt}
  \renewcommand{\arraystretch}{1.35}
  \begin{tabularx}{\linewidth}{@{}lX@{}}
    \toprule
    \textbf{Method} & \textbf{Selected Features} \\
    \midrule
    Lasso (L1)    & EF, SC, CPK, age \\
    Elastic Net   & EF, SC, age, sodium, CPK\ldots \\
    \textbf{MRMR} & \textbf{SC + EF only} \\
    \bottomrule
  \end{tabularx}

  \medskip
  \textbf{Random Forest importance scores:}

  \medskip
  \foreach \feat/\val/\w in {
    {Serum Creatinine}/{0.227}/9.5,
    {Ejection Fraction}/{0.189}/7.9,
    {CPK}/{0.138}/5.8,
    {Platelets}/{0.125}/5.2,
    {Age}/{0.113}/4.7,
    {Serum Sodium}/{0.100}/4.2%
  }{%
    \begin{tikzpicture}[baseline=-0.07cm]
      \node[anchor=west, font=\normalsize, minimum width=5.2cm, text=TGray]
        at (0,0) {\feat};
      \fill[Blue!55] (5.5cm,-0.14cm) rectangle (5.5cm+\w cm,0.30cm);
      \node[anchor=west, font=\normalsize\bfseries, text=Navy]
        at (5.5cm+\w cm+0.12cm,0.08cm) {\val};
    \end{tikzpicture}\\[2pt]
  }

  \medskip
  \begin{tcolorbox}[findcard=Green]
    \textbf{\color{Green}Key result:}
    Using only EF + SC, RF achieves
    \textbf{Acc = 78.9\%, F1 = 0.655} —
    matching full-feature models.
  \end{tcolorbox}
\end{block}

\end{column}

%=============================================================
% COL 3 — Model Performance chart + table
%=============================================================
\begin{column}{28.8cm}

\begin{block}{Model Performance — Test Set (n\,=\,90)}

  \medskip
  % Pure TikZ horizontal bar chart — guaranteed label display
  \begin{tikzpicture}[x=1cm, y=1.0cm]
    \foreach \lbl/\val/\clr/\i in {
      MLP/0.678/gray!40/0,
      {PLS-DA}/0.722/gray!45/1,
      KNN/0.667/gray!40/2,
      {Log.\ Reg.}/0.711/Blue!35/3,
      {Grad.\ Boost.}/0.733/Blue!50/4,
      {SVM (RBF)}/0.733/Blue!65/5,
      {RF Baseline}/0.733/Blue!78/6,
      {2-Feat.\ RF}/0.789/Amber!80/7,
      {LightGBM (T)}/0.778/Green!78/8%
    }{
      \pgfmathsetmacro{\blen}{(\val-0.60)/0.24*15}
      \fill[\clr,rounded corners=2pt]
        (0,\i+0.08) rectangle (\blen,\i+0.88);
      \node[anchor=east, font=\normalsize, text=TGray]
        at (-0.15,\i+0.48) {\lbl};
      \node[anchor=west, font=\small\bfseries, text=TGray]
        at (\blen+0.12,\i+0.48) {\val};
    }
    % grid lines + tick labels
    \foreach \v/\lb in {0.65/{65\%},0.70/{70\%},0.75/{75\%},0.80/{80\%}}{
      \pgfmathsetmacro{\tx}{(\v-0.60)/0.24*15}
      \draw[gray!25,dotted,thin] (\tx,-0.05)--(\tx,9.05);
      \node[anchor=north,font=\small,text=LGray] at (\tx,-0.1) {\lb};
    }
    \draw[gray!40,thin] (0,-0.05)--(15,-0.05);
    \node[font=\normalsize,text=LGray] at (7.5,-0.7)
      {Accuracy on held-out test set};
  \end{tikzpicture}

  \medskip
  \setlength{\tabcolsep}{5pt}
  \renewcommand{\arraystretch}{1.35}
  \begin{tabularx}{\linewidth}{@{}Xcccc@{}}
    \toprule
    \textbf{Model} & \textbf{Acc} & \textbf{F1} & \textbf{MCC} & \textbf{AUC} \\
    \midrule
    \rowcolor{LGreen}
    \textbf{LightGBM}\,$^\star$ & \textbf{.778} & \textbf{.667} & \textbf{.501} & .747 \\
    \rowcolor{LAmber}
    2-Feat.\ RF\,$^\dagger$     & .789 & .655 & —    & —    \\
    RF Baseline                 & .733 & .600 & .401 & .760 \\
    Grad.\ Boosting             & .733 & .586 & .390 & .709 \\
    SVM (RBF)                   & .733 & .520 & —    & \textbf{.768} \\
    Logistic Reg.               & .711 & .458 & —    & .746 \\
    KNN                         & .667 & .348 & —    & .643 \\
    PLS-DA                      & .722 & .444 & —    & —    \\
    MLP (500\,ep.)              & .68  & —    & —    & —    \\
    \bottomrule
  \end{tabularx}
  \smallskip
  {\small $^\star$\,Optuna-tuned.\enspace
          $^\dagger$\,EF\,+\,SC only.\enspace
          MCC accounts for class imbalance.}
\end{block}

\end{column}

%=============================================================
% COL 4 — Deep Learning · XAI · Key Findings · References
%=============================================================
\begin{column}{28.8cm}

\begin{block}{Deep Learning — Week 8}
  \textbf{PyTorch MLP architecture:}
  Input(11) $\to$ Dense(32, ReLU) $\to$ Sigmoid(1)

  \medskip
  Loss: Binary cross-entropy \enspace|\enspace
  Optimiser: Adam \enspace|\enspace
  Epochs: 500 \enspace|\enspace
  Early stopping on val loss

  \medskip
  \setlength{\tabcolsep}{6pt}
  \renewcommand{\arraystretch}{1.35}
  \begin{tabularx}{\linewidth}{@{}XXXX@{}}
    \toprule
    \textbf{Epoch} & \textbf{Train Loss} & \textbf{Val Loss} & \textbf{Test Acc} \\
    \midrule
    50  & 0.549 & 0.553 & 68.9\% \\
    100 & 0.510 & 0.527 & 67.8\% \\
    200 & 0.480 & 0.490 & 67.8\% \\
    \bottomrule
  \end{tabularx}

  \medskip
  Deep learning \emph{underperforms} tree ensembles on this
  $n$\,=\,299 tabular dataset — consistent with the broader
  literature on structured clinical data.
\end{block}

\begin{block}{Explainability: SHAP \& PLS-DA — Week 9}
  \textbf{SHAP} applied to Random Forest and MLP
  (KernelExplainer, 100 background samples, all 299 patients
  explained):

  \medskip
  \begin{itemize}
    \item \textbf{Global (beeswarm):} Serum creatinine and
          ejection fraction dominate mean $|\text{SHAP}|$
          across all patients
    \item \textbf{Local (waterfall):} Elevated SC pushes
          predictions toward death; high EF pushes toward
          survival — interpretable per patient
    \item \textbf{Model-agnostic:} RF and MLP agree on the
          same top features, confirming result robustness
  \end{itemize}

  \medskip
  \textbf{PLS-DA:} 2 latent variables achieve Acc = 72.2\%.
  Most discriminative signal is concentrated in a subspace
  aligned with ejection fraction and serum creatinine.
\end{block}

\begin{block}{Key Findings}
  \begin{tcolorbox}[findcard=Green]
    \textbf{\color{Green}2 Features Sufficient} —
    EF + SC alone: Acc 78.9\%, F1 0.655
  \end{tcolorbox}
  \vspace{4pt}
  \begin{tcolorbox}[findcard=Blue]
    \textbf{\color{Blue}Best Overall: LightGBM (Tuned)} —
    Acc 77.8\%, F1 0.667, MCC 0.501
  \end{tcolorbox}
  \vspace{4pt}
  \begin{tcolorbox}[findcard=Amber]
    \textbf{\color{Amber}Best Ranking: SVM (RBF)} —
    AUC-ROC = 0.768
  \end{tcolorbox}
  \vspace{4pt}
  \begin{tcolorbox}[findcard=Navy]
    \textbf{\color{Navy}SHAP Consistent} —
    RF and MLP agree on feature rankings
  \end{tcolorbox}
\end{block}

\begin{block}{References}
  Chicco \& Jurman (2020). \textit{Machine learning can predict
  survival of patients with heart failure from serum creatinine
  and ejection fraction alone.}
  \textit{BMC Medical Informatics and Decision Making}, 20(16).

  \medskip
  UCI ML Repository: Heart Failure Clinical Records.
  \textbf{Tools:} Python, scikit-learn, LightGBM, Optuna,
  PyTorch, SHAP.

  \medskip
  \textbf{GitHub:} MichiganDataScienceTeam/
  W26-MDST-Project\_Heart-Failure-Survival-Analysis
\end{block}

\end{column}
\end{columns}

\end{frame}
\end{document}
